\section{Proof assurance and verifier trust}

Formal verification sharply reduces one class of error but introduces a trust boundary. Lean uses a small kernel to check proof terms; tactic bugs do not normally threaten kernel soundness because tactics must produce terms accepted by the kernel \cite{leanref}. Finished developments must nevertheless audit their transitive axioms. Lean's documentation states that \texttt{sorryAx}, used by \texttt{sorry}, can prove arbitrary propositions and is not intended in finished proofs \cite{leanaxioms}.

The strict profile records a proof receipt containing theorem and source hashes, checker identity and version, transitive axioms, and an independent recheck where supported. It rejects \texttt{sorryAx}, unregistered custom axioms and, when independent kernel-level assurance is required, proof paths that depend on compiler/native trust. This last condition is motivated by the fact that formal verification itself has an implementation trust surface: Lean 4.32.1, released in July 2026, fixed a kernel soundness bug and noted that the recommended separate \texttt{comparator} path for potentially dishonest proofs was not affected \cite{lean4321}.

\begin{proposition}[Generator-error containment]
Let $G$ be any proposal generator, $K$ a sound proof checker for logic $L$, and $U$ the canonical update rule. Suppose
\[
U(T,p)=\mathrm{promote}\Longrightarrow K(p,T)=\mathrm{accept}.
\]
Then, relative to the soundness of $K$ and the registered axioms, an invalid theorem cannot enter canonical truth state solely because $G$ generated a plausible but invalid proof step.
\end{proposition}

\begin{proof}
The generator has no canonical write authority. Every promoted theorem must have an accepted proof artifact. By soundness of $K$, acceptance implies derivability of $T$ from the registered assumptions. Therefore an invalid proposal generated by $G$ is either rejected or would require failure of a trust assumption outside the generator itself. \qed
\end{proof}

\begin{proposition}[Verified checkpointing converts local generator error into search cost]
Assume a final theorem depends only on lemmas admitted to a persistent DAG after independent proof checking. Then the per-proposal probability that an LLM suggests an invalid local step does not multiply directly into the logical validity of the accepted theorem. Invalid proposals instead increase rejected branches, latency or compute, unless a specification or verifier-trust assumption fails.
\end{proposition}

\begin{proof}
An invalid proposal that fails checking is not added to the trusted dependency DAG, hence cannot serve as a premise of the accepted theorem. The final theorem's logical status depends on the accepted proof terms and checker assumptions, not on the number of invalid discarded proposals. \qed
\end{proof}

This proposition does not make research easy. It changes the failure geometry. Long-horizon model error remains costly, and the hard unsolved problems become search, formalization, representation, novelty and trust rather than hidden survival of one invalid implication.
